\documentclass[10pt]{article}
\usepackage[utf8]{inputenc}
\usepackage[T1]{fontenc}
\usepackage{amsmath}
\usepackage{amsfonts}
\usepackage{amssymb}
\usepackage[version=4]{mhchem}
\usepackage{stmaryrd}

\setlength{\parindent}{0pt}
\setlength{\parskip}{6pt}

% Space between parts of a question
\newcommand{\qspace}{\vspace{5cm}}
% Space between full questions
\newcommand{\bigqspace}{\vspace{4cm}}

\begin{document}
\textbf{MATH 261, Homework 10} \hfill Name: \underline{\hspace{6cm}}

\bigskip

%-----------------------------------------------------------------------


\begin{enumerate}
  \item Compute the divergence of the vector field $\mathbf{F}=\sin (x+z) \mathbf{i}-y e^{x z} \mathbf{k}$
	
	 \vspace*{0.35\textheight}
	
  \item Find a constant $c$ for which the velocity field $\mathbf{v}=(c x-y) i+(y-z) j+(3 x+4 c z) k$ of a fluid is incompressible (i.e., $\operatorname{div}(\mathbf{v})=0$).
	
	 \vspace*{0.5\textheight}
	
  \item Let $S$ be the outward oriented sphere $x^{2}+y^{2}+z^{2}=1$. Let $\mathbf{F}=\left\langle\frac{1}{2} x^{2}, 3,5\right\rangle$ be a vector field. Use divergence theorem to compute the flux integral $\iint_{S} \mathbf{F} \cdot \mathbf{N} d S$.
	 
	\vspace*{0.42\textheight}
	
  \item Let $S$ be the outward oriented cube bounded by the planes $x= \pm 1, y= \pm 1$, and $z= \pm 1$. Let $\mathbf{F}=\left\langle x y^{2}, x^{2}, z^{2}\right\rangle$ be a vector field. Use divergence theorem to compute the flux integral $\iint_{S} \mathbf{F} \cdot \mathbf{N} d S$.
	 \vspace*{0.41\textheight}
	
  \item Let $S$ be the outward surface consisting of the portion of the sphere given by $x^{2}+y^{2}+ z^{2}=4$ between the planes $z=-1$ and $z=1$, as well as the portion of the plane $z=1$ and $z=-1$ which is within the sphere (so the surface is closed). Let $\mathbf{F}=\left\langle x, y, x^{2}\right\rangle$ be a vector field. Use divergence theorem to compute the flux integral $\iint_{S} \mathbf{F} \cdot \mathbf{N} d S$.
	
	 \vspace*{0.35\textheight}
	
  \item Let $S$ be the outward surface consisting of the portion of the cone $z=\sqrt{x^{2}+y^{2}}$ which is above the plane $z=1$ and below the plane $z=2$, along with the portions of the planes $z=1$ and $z=2$ which are within the cone, so the surface is closed. Let $\mathbf{F}=\left\langle e^{y}, 5, \frac{1}{2} z^{2}\right\rangle$ be a vector field. Use divergence theorem to compute the flux. integral $\iint_{S} \mathbf{F} \cdot \mathbf{N} d S$.
  
	\vspace*{0.4\textheight}
	
	\item Let $S$ be the outward surface consisting of the portion of the cylinder $x^{2}+y^{2}=9$ which lies above the plane $z=0$ and below the plane $3 x+2 y+5 z=120$, as well as the portions of both planes which lie inside the cylinder(so that the surface is closed). Let $\mathbf{F}=\langle 5,3,6\rangle$ be a vector field. Use divergence theorem to compute the flux integral $\iint_{S} \mathbf{F} \cdot \mathbf{N} d S$.
  
	 \vspace*{0.33\textheight}
	
	\item Let $S$ be the outward oriented surface consisting of the portion of cone $z=\sqrt{x^{2}+y^{2}}$ which is below the sphere $x^{2}+y^{2}+z^{2}=9$ and above the plane $z=1$, as well as the portions of the sphere $x^{2}+y^{2}+z^{2}=9$ and plane $z=1$ which are within the cone (so the surface is closed). Let $\mathbf{F}=\left\langle x y,-\frac{1}{2} y^{2}, z\right\rangle$ be a vector field. Use divergence theorem to compute the flux integral $\iint_{S} \mathbf{F} \cdot \mathbf{N} d S$.
  
	 \vspace*{0.4\textheight}
	
	\item Let $S$ be the outward surface consisting of the cylinder $x^{2}+y^{2}=1$ and $x^{2}+y^{2}=4$ for $0 \leq z \leq 1$, as well as the portions of both planes $z=0$ and $z=1$ which are between the two cylinders (so that the surface is closed). Let $\mathbf{F}=\left\langle x y^{2}, y x^{2}, x^{2} y^{2}\right\rangle$ be a vector field. Use divergence theorem to compute the flux integral $\iint_{S} \mathbf{F} \cdot \mathbf{N} d S$.
	
	 \vspace*{0.33\textheight}
  \item Let $S$ be the outward surface consisting of the portion of the cylinder $x^{2}+y^{2}=4$ which is between the plane $z=0$ and the paraboloid $z=x^{2}+y^{2}$, as well as the portion of the plane and the paraboloid which are within the cylinder (so that the surface is closed). Let $\mathbf{F}=\langle y, x y,-z\rangle$ be a vector field. Use divergence theorem to compute the flux integral $\iint_{S} \mathbf{F} \cdot \mathbf{N} d S$.
	
	\vspace*{0.4\textheight}
	
  \item Let $S$ be the outward surface consisting of the portion of the cylinder $x^{2}+y^{2}=4$ which is within the sphere $x^{2}+y^{2}+z^{2}=20$, as well as two portions of the sphere (so that the surface is closed). Let $\mathbf{F}=\langle x z,-y z, z x\rangle$ be a vector field. Use divergence theorem to compute the flux integral $\iint_{S} \mathbf{F} \cdot \mathbf{N} d S$.
  
	\vspace*{0.4\textheight}
	
	\item Find the maximum value of the function $f(x, y)=x^{2}+y$ along the curve $x^{2}+y^{2}=1$ using the method of Lagrange multipliers.
	
	\vspace*{0.42\textheight}
	
  \item Find the minimum value of the function $f(x, y)=x^{2}+y^{2}$ along the curve $x y=1$ using the method of Lagrange multipliers. At what point (points) does it occur?
	
	\vspace*{0.42\textheight}
  
	\item Find the maximum value of the function $f(x, y, z)=y z+x y$ along the curve $x^{2}+ y^{2}+z^{2}=1$ using the method of Lagrange multipliers. At what point (points) does it occur?
	
	\vspace*{0.42\textheight}
	
  \item Find the maximum value of the function $f(x, y, z)=z e^{x y}$ along the curve $x^{2}+y^{2}+z^{2}=$ 1 using the method of Lagrange multipliers. At what point (points) does it occur?
  
	\vspace*{0.42\textheight}
	
	\item Find the extreme values of $f(x, y)=x y$ along the curve $x^{2}+y^{2}=10$ using the Lagrange multiplier method.
   
	\vspace*{0.45\textheight}
	
	\item Find the maximum value of the function $f(x, y)=49-x^{2}-y^{2}$ on the line $x+3 y=10$ using the method of Lagrange multipliers. At what point (points) does it occur?
  
	 \vspace*{0.42\textheight}
	
	\item Find the points on the curve $x y^{2}=54$ closest to the origin.
  
	 \vspace*{0.46\textheight}
	
	\item Find the dimensions of the closed right circular cylindrical can with the smallest surface area whose volume is $16 \pi \mathrm{~cm}^{3}$.
	
	 \vspace*{0.42\textheight}
	
  \item Set up, but do not solve, the system of equations you would need to solve to find the extreme value of the function $f(x, y, z)=x^{2} y+z^{2} y+x^{2} y^{2}$ along the curve $x^{4}+y^{4}+z^{4}=1$.
	
	 \vspace*{0.45\textheight}
	
  \item Compute the first-degree Taylor approximation to the function $f(x, y)=x^{2} y+x^{3}+1$ at the point $(1,1)$.
	
	 \vspace*{0.15\textheight}
	
  \item Compute the first-degree Taylor approximation to the function $f(x, y)=y^{4} e^{2 x}$ at the point $(0,1)$.
  
	\vspace*{0.2\textheight}
	
	\item Compute the first-degree Taylor approximation to the function $f(x, y)=\sin (x) \cos (y)$ at the point $(0,0)$.
  
	\vspace*{0.2\textheight}
	
	\item Compute the first-degree Taylor approximation to the function $f(x, y)=\sqrt{6-x^{2}-y^{2}}$ at the point $(1,1)$.
		
		\vspace*{0.2\textheight}
	
  \item Compute the second-degree Taylor approximation to the function $f(x, y)=x^{3} e^{y}$ at the point $(1,0)$.
  
		\vspace*{0.25\textheight}
		
	\item Compute the first-degree Taylor approximation to the function $f(x, y)=x^{2} y^{2}+e^{y^{2}}$ at the point $(1,1)$.
  
		\vspace*{0.25\textheight}
	
	\item Suppose $f(x, y)$ is a function with $f(0,0)=5, f_{x}(0,0)=1, f_{y}(0,0)=0, f_{x x}(0,0)=$ -3 , $f_{y y}(0,0)=0$ and $f_{x y}(0,0)=3$. Compute the second-degree Taylor approximation.
\end{enumerate}

\end{document}